\chapter*{Preface}
\addcontentsline{toc}{chapter}{Preface}

Formal logic studies formal languages, interpretations of those languages,
and rule-governed derivations.  Correct work in formal logic therefore requires
more than symbol manipulation.  It requires exact control of the level of
language, the grammar of expressions, variable binding, interpretation,
assumptions, inference rules, and the metatheoretic hypotheses under which a
result is stated.

This reference is organized at graduate level rather than as a catalogue of
disconnected systems.  The dependency order is deliberate.  Syntax is
defined before semantics; semantics is defined before soundness; a calculus is
defined before derivability; classical propositional logic is established
before first-order logic; classical first-order logic is established before
model theory, computability, incompleteness, proof theory, and nonclassical
logic.  The organization follows the core treatments in
\cite{Enderton2001,EbbinghausFlumThomas1994,OpenLogic2026,Marker2024}, while the
historical and advanced results are tied to their original sources where
practical.

The document has four purposes:
\begin{enumerate}
  \item to teach the formal-logical content of the subject;
  \item to teach the construction and checking of derivations, proofs,
        countervaluations, structures, and countermodels of sentence
        arguments;
  \item to teach the writing and typesetting of formal logic; and
  \item to provide a dependency-ordered reference for graduate study, teaching,
        and preparation for qualifying examinations.
\end{enumerate}

The reference is self-contained at the level of notation and foundational
definitions.  Advanced model theory and proof theory are presented to the
point required for a broad graduate-level reference, but a research specialization
still requires the monographs cited in the relevant chapters.

